\section{虚树}

当若干询问各自只涉及规模为 $k$（$k\ll n$）的点集 $S$、且要在树上按父子关系做 DP 时，抽取 $S$ 及其两两 LCA 组成虚树，边数不超过 $2k-1$，在虚树上做同样转移即可。

需已有以 $1$ 为根的原树：\texttt{n} 为点数，\texttt{dfn[u]} 为 DFS 进入序（保证子树连续即可），\texttt{dep[u]} 为深度，并可用 \verb|lca(u,v)| 求最近公共祖先（即本专题「最近公共祖先」任一版本给出的量）。\verb|buildVirtualTree(S)| 返回虚树邻接表 \texttt{vt}（无向存边），根为 $S$ 中点的 LCA；未涉及点为空。每次调用新建邻接表，可重复调用。

排序与建树复杂度 $\mathcal{O}(k\log k)$，另含 $\mathcal{O}(k)$ 次 LCA 查询。

\begin{lstlisting}
auto buildVirtualTree = [&](std::vector<int> S) -> std::vector<std::vector<int>> {
    std::sort(S.begin(), S.end(), [&](const auto& x, const auto& y) {
        return dfn[x] < dfn[y];
    });
    S.erase(std::unique(S.begin(), S.end()), S.end());
    std::vector<std::vector<int>> vt(n + 1);
    std::vector<int> stk;
    auto link = [&](int p, int c){
        vt[p].push_back(c);
        vt[c].push_back(p);
    };
    for (int u : S){
        if (stk.empty()){ stk.push_back(u); continue; }
        int w = lca(stk.back(), u);
        while ((int)stk.size() >= 2 && dep[w] <= dep[stk[stk.size() - 2]]){
            link(stk[stk.size() - 2], stk.back());
            stk.pop_back();
        }
        if (stk.back() != w){
            link(w, stk.back());
            stk.pop_back();
            if (stk.empty() || stk.back() != w) stk.push_back(w);
        }
        stk.push_back(u);
    }
    while ((int)stk.size() >= 2){
        link(stk[stk.size() - 2], stk.back());
        stk.pop_back();
    }
    return vt;
};
\end{lstlisting}
